%!TEX root = ../../main.tex
\chapter{A 业务建模}
% ============================================================

\section{案例背景}

老板（柜台）在商场经营一个玉器门店（档口），在柜 4–5 人，销售手镯、吊坠、摆件等成品玉器。
玉器经营与普通零售不同：每件玉器都是独一件、一物一价，客单价从一两千元到六位数不等；
在柜货品约 3000–4000 件（小件多、部分货品可拆，数量常有出入），货值约 2000–3000 万元，
其中相当一部分是向同行调来的货，并非都是门店自有。
货品来源按占比大致是：同行调货约 40\%、买料加工约 30\%、回收回流约 20\%、成品进货约 10\%；
主做新疆料，加工以切山料、籽料为主，戈壁料也卖一些，原石山流水料因风险大而少买。
正因为非标准化——同一款手串做 100 条，克重、成色、定价各不相同，无法按“件”套用现成进销存，
只能逐件建档。玉的真假依靠鉴定证书，而产地目前证书做不出来，靠老板眼力把关。

随着经营规模扩大，管理方式没有相应跟上，问题集中在四处：

\begin{enumerate}[leftmargin=1.6em, itemsep=2pt]
  \item 调货往来靠微信：调货时在图片上写价、发图给同行，货拿走就算成交，
        微信记录会丢、图片会被清理，事后认不出调出去的是哪一件，对账经常扯皮——这是最头疼的一处；
  \item 进销记录靠纸质台账和记忆：每件货品从入库到售出没有系统记录，
        成本和入库时间要翻进货单据回忆，货品周转周期算不出来，压柜的货越来越多；
  \item 客户信息散落：客户联系方式、购买记录、定金与改圈口等约定散落在微信聊天和手写小票里，
        想做精准营销没有数据可用，月度销售客户数要月底翻账本数出来；
  \item 经营分析缺数据：销售额靠月底翻账本数出来，按货品、客群的分析无法进行。
\end{enumerate}

为此决定开发玉器门店顾客服务与经营管理系统，把每件货品从进货入库到销售回款的全程记录起来。

上述信息来自对门店老板的当面访谈。2026 年 10 月 1 日在门店访谈一次，围绕门店规模、
记账方式、经营痛点、货源、定价与售后逐项提问并现场记录；另以公众号“长生玉鉴”的
三篇文章作背景资料，补充行业的通行做法。

选访谈而不做问卷，是因为口径只集中在老板一人手里：在柜人数、货值、货源占比、毛利，
柜台销售答不上来。玉器又非标，同一款手串做一百条，克重、成色、定价各不相同，
只有当面追问才能问清“为什么这样定价”“调货什么时候算成交”这类细节。
访谈的局限是单一视角：老板讲的是经营结果，具体单据的格式与栏目覆盖不到，
因此涉及小票、进货单的部分，以门店现用的纸质单据实物为准，不凭访谈推断。

\section{涉众与期望}

\needspace{8\baselineskip}
\begin{table}[H]
\centering
\caption{涉众及其利益}
\begin{tabularx}{\textwidth}{@{}l l Y@{}}
\toprule
涉众 & 与业务的关系 & 涉众利益 \\
\midrule
门店老板 & 业务工人、经营决策者 & 掌握每件货品的进销存与周转周期，减少积压；月度销售额与客户数随手可查；回收、调货往来不错账，出错了追得到经手人 \\
销售 & 业务工人、一线执行者 & 记录过程简单、不干扰成交；成交归属清楚；老板不在店时代理收货与送检也有据可查 \\
客户 & 业务执行者 & 交易顺畅，价格清楚，证书可信，售后（退货）有据，鉴赏期内能从容决定 \\
玉料供应商、成品玉器供应商 & 业务执行者 & 供货与结算账目清楚 \\
加工师傅 & 业务执行者 & 加工要求、交期与加工费明确，加工过程可回看 \\
鉴定机构 & 业务执行者 & 送检样品与证书信息完整 \\
回流出让方 & 业务执行者 & 回收作价合理、付款及时 \\
同行商户 & 业务执行者 & 调货往来账目清楚，不再靠微信图片对账 \\
\bottomrule
\end{tabularx}
\end{table}

\noindent
用例规约中的每条约束都回溯到上表中某条利益，见 B 部分各用例规约的涉众利益与补充约束栏。

\section{愿景}

愿景写的是引进系统后目标组织要改进的指标。

\needspace{8\baselineskip}
\begin{table}[H]
\centering
\caption{愿景（目标系统：玉器门店顾客服务与经营管理系统）}
\begin{tabularx}{\textwidth}{@{}L{2.6cm} Y@{}}
\toprule
项目 & 内容 \\
\midrule
系统名称 & 玉器门店顾客服务与经营管理系统 \\
目标组织 & 玉器销售门店柜台（档口老板 + 销售） \\
负责人 & 门店老板（柜台） \\
业务目标 & 提高门店销售额，降低货品周转周期 \\
度量指标 & ① 月度销售客户数；② 货品加权周转周期（统计口径：从货品进货入库开始，至货品销售完成并收回货款为止，按货品总金额加权平均） \\
对比时间 & 2025 年全年业务月均值（引入系统前，沿用纸质台账的做法） \\
改进值 & 销售额提升 10\%；货品周转周期下降 20\% \\
\bottomrule
\end{tabularx}
\end{table}

\begin{itemize}[leftmargin=1.6em, itemsep=2pt]
  \item 玉器一物一价、金额差异大，周转周期必须按件追踪：从进货入库开始，
        到销售完成并收回货款为止，按货品总金额加权平均。六位数的货压三四年是常事，
        几百到一两千的货一周多就能走，不加权的话平均值会掩盖高价货的真实积压。
        月度销售客户数按当月成交客户去重计。
  \item 改进值是拟议目标，不是承诺；周转周期下降 20\% 来自高价压柜货品及早发现和调剂，
        以系统记录为统计依据。销售额提升 10\% 主要来自老客户复购与线上引流到店的转化——
        收料成本和定价靠眼力，系统不决定成交价，只把货品与客户记录支撑起来。
\end{itemize}

\section{业务用例图}

业务用例图描述目标组织对外提供的价值。目标组织是玉器门店（档口），
业务执行者共七位：客户、回流出让方、玉料供应商、加工师傅、成品玉器供应商、鉴定机构、同行商户。
业务用例图见图 \ref{fig:01-business-usecase.png}；手绘稿见图 \ref{fig:01a-business-usecase-hand.jpg}。

\figpab[0.86]{业务用例图（手绘稿）}{01a-business-usecase-hand.jpg}

\figab[0.9]{业务用例图（PlantUML 绘制）}{01-business-usecase.png}

\noindent
图中组织边界内共有 11 个业务用例，名称按业务执行者提出请求的口径命名：
面向客户侧的定制、购买玉器、退货，
面向货品来源侧的回收（收购回流玉器）、出售玉料、接加工、出售成品玉器、送检、调货，
以及承载货款结算的支取工费、支付供应商货款。
货品来源侧的四个方向（出售玉料、出售成品玉器、回收、调货）都对应同一条货款结算线，
所以支付供应商货款画成一个用例、由多方共同参与。
定制由门店与客户线下谈定料子、器型与交期，不经过系统。

引入目标系统后，新增了面向客户的线上服务用例（预定商品、预约看货）、面向合作商户的查看调货价目、面向供应商的提交出让申请与面向加工师傅的查看加工任务，改进业务用例图见图 \ref{fig:02-business-usecase-to-be.png}。

\figab[0.9]{改进业务用例图：新增的线上服务用例}{02-business-usecase-to-be.png}

\section{现状业务序列图}

现状业务序列图描述组织当前的实际做法。
图中人员与记录工具都照实画（纸质销售台账、进货单据都是业务实体），
消息统一写成“A 请求 B 做某事”的服务调用，名字用“动词+宾语()”，
代表的是责任与目的，跟数据流向无关：箭头指向谁，就由谁承担这项工作，
完成后用虚线把结果返回给请求方，处理期间画激活条。
符号约定：业务执行者（客户）画小人；业务工人（销售、老板）画圆圈内人形，
圆圈右下加一道斜线切口、顶部带小箭头；业务实体（纸质销售台账、进货单据）画圆圈加右下斜线，
圆圈底下加一横；对象名统一写成加粗下划线的“:名字”。

现状的日常经营拆成场景逐一定格：现场议价成交、退货、成品进货、
同行调货（调入、调出各一张）、玉料采购与加工。
这些流程里，问题出在四处：
调货靠微信在图片上写价、发图成交，记录会丢、图片会被清理，事后认不出调出去的是哪件货，
对账只能靠回忆；售出只在台账记一笔金额，没有和具体货品关联，货品从入库到售出的时间无从对齐，
周转周期只能估算；客户成交后不留任何记录，月度销售客户数只能翻账本数；
月底汇总靠人工，忙时漏记当晚无从发现。

\figpab[0.9]{现状场景：议价成交}{02-as-is-sale.png}
\figpab[0.9]{现状场景：退货}{03-as-is-return.png}
\figpab[0.9]{现状场景：成品进货}{04-as-is-goods-in.png}
\figpab[0.9]{现状场景：同行调货（合作商户调货给我们）}{05-as-is-transfer-in.png}
\figpab[0.9]{现状场景：同行调货（调出货品给合作商户）}{06-as-is-transfer-out.png}
\figab[0.9]{现状场景：玉料采购与加工}{07-as-is-material.png}

\section{改进业务序列图}

改进思路是把人工记账与统计交给新引入的系统：新系统是新的业务实体，
接手纸质销售台账与进货单据的部分责任。货品入库时逐件建档（编号、品名、来源、成本、入库时间），
成交时按货品编号登记并更新货品状态，销售与周转统计由系统按件自动累计；
保存、校验、统计都是系统内部行为，画成自消息，不单独画数据库生命线。

六个现状场景逐一改进：凡是改进场景里从外部指向目标系统的消息，
都映射为图~\ref{fig:08-system-usecase.png} 的系统用例——登记成交、登记退货、登记货品入库、
预定商品、预约查看货品（图中消息为“预约看货”）、查看调货价格（消息为“查看调货价目”）、
提交货品出让申请、审核出让申请、预约上门时间、查看加工任务、提交完工交付申请，
以及由各货款消息聚合的登记货款记录、由调货往来与库存消息聚合的维护调货记录，共 13 个；
维护客户信息、登记送检与证书、查看销售与周转统计、维护基础数据、登录、浏览货品、
开通外部账户 7 个用例是支撑上述用例与愿景指标所必需的补充，在图中一并列出。

六个现状场景的改进版逐一如下：线上预定后到店成交、退货登记与货品状态回柜、
成品进货（供应商在线提交出让申请并预约上门）、同行调货（调入、调出）、
玉料采购与加工（加工师傅在线查看任务与提交完工交付）。
六个现状场景之外补一个新增业务场景——供应商寄售（改进场景图~\ref{fig:15-to-be-consignment.png}），
该场景原来不存在，没有现状图，直接给出改进版：
寄售货品凭合同建档入库，售出后按寄售价结算，到期未售出货品退回。

\figpab[0.9]{改进场景：议价成交（线上预定 + 成交当场登记）}{09-to-be-sale.png}
\figpab[0.9]{改进场景：退货}{10-to-be-return.png}
\figpab[0.9]{改进场景：成品进货（在线出让申请与预约上门）}{11-to-be-goods-in.png}
\figpab[0.9]{改进场景：同行调货（合作商户调货给我们）}{12-to-be-transfer-in.png}
\figpab[0.9]{改进场景：同行调货（调出货品给合作商户）}{13-to-be-transfer-out.png}
\figpab[0.9]{改进场景：玉料采购与加工（在线查看任务与交付）}{14-to-be-material.png}
\figpab[0.9]{改进场景：供应商寄售（合同入库与寄售结算）}{15-to-be-consignment.png}

\section{业务改进对照}

\needspace{8\baselineskip}
\begin{table}[H]
\centering
\caption{现状痛点与改进后对照}
\begin{tabularx}{\textwidth}{@{}L{3.2cm} Y Y@{}}
\toprule
环节 & 现状 & 改进后 \\
\midrule
调货往来 & 微信发图写价，记录会丢，对账扯皮 & 调货往来与调货库存进系统，按货品编号与克重对账 \\
成交记录 & 手写小票 + 台账，只有金额 & 按货品编号登记成交，必填货品与成交价 \\
货品档案 & 无（进货信息靠单据和记忆） & 逐件建档：编号、品名、来源、成本、入库时间 \\
周转统计 & 无法统计 & 系统按件记录入库与售出回款，按金额加权计算周转周期 \\
客户与营销 & 散落在微信和记忆里 & 客户档案关联购买记录，支撑精准营销 \\
月度汇总 & 月底人工翻账本 & 自动汇总月度销售额与销售客户数 \\
\bottomrule
\end{tabularx}
\end{table}

六处改进可归为三类。调货往来、成交记录、货品档案三行属于把物流变成信息流：
货品流转与价格原来靠微信图片和纸质小票，改进后逐件进系统，调出哪一件、什么价、
谁经手都有记录。周转统计、客户与营销、月度汇总三行属于信息流自动化：原来看不出周转、
客户与购买记录对不上、月底要翻账本，改进后由系统按入库与回款时间自动计算，
客户档案与成交记录自动关联。底价校验、货品状态按状态机迁移、定金与鉴赏期决定结算时点
这类规则，集中在业务层一处，柜台端与小程序共用同一套判断，属于封装领域逻辑——
规则改动只改一处，两个前端不会给出不同结论。

% ============================================================
